% !PRINTED 345
\weberpage{345}

Dieser Ausdruck gilt aber allgemein. Denn sind $m_1$ und $m_2$
relativ prim und ist
\[
m_1 = N(\mathfrak{m}_1), \quad m_2 = N(\mathfrak{m}_2),
\]
so ist
\[
m = m_1 m_2 = N(\mathfrak{m}_1 \mathfrak{m}_2).
\]

Man kann also aus einer Darstellung von $m_1$, $m_2$ als Normen
eine Darstellung von $m$ als Norm ableiten.

Ist umgekehrt
\[
m_1 m_2 = N(\mathfrak{m}) = \mathfrak{m}\,\mathfrak{m}',
\]
wo $\mathfrak{m}$, $\mathfrak{m}'$ konjugierte Ideale sind, so muß, wenn $m_1$ und $m$ durch
ein Primideal $\mathfrak{p}$ teilbar sind, $m_1$ auch durch $\mathfrak{p}\,\mathfrak{p}'$ teilbar sein, also
ist auch $m_1 = N(\mathfrak{m}_1)$ und ebenso $m_2 = N(\mathfrak{m}_2)$.

Demnach gilt, wenn $m_1$ und $m_2$ relativ prim sind:
\weq{7}{\psi(m_1)\,\psi(m_2) = \psi(m_1 m_2).}
Andererseits ist, wenn $n_1$ und $n_2$ die Teiler von $m_1$ und $m_2$
durchlaufen,
\[
\Sigma(\Delta,\, n_1)\, \Sigma(\Delta,\, n_2) = \Sigma(\Delta,\, n_1 n_2),
\]
und damit ist die Formel (6) allgemein nachgewiesen. Wir
sprechen dies als Satz aus:

4. \gesperrt{Eine natürliche Zahl $m$ kann in dem quadratischen
Körper $\Omega$ mit der Grundzahl $\Delta$ auf}
\[
\Sigma(\Delta,\, n)
\]
\gesperrt{verschiedene Arten als Norm eines Ideals dargestellt werden,
wenn $n$ die sämtlichen Teiler von $m$ durchläuft.}

\S\,94. \gesperrt{Das quadratische Reziprozitätsgesetz.}

Aus der Zerlegung der Primzahlen im quadratischen Körper
leitet \gesperrt{Dedekind} einen neuen Beweis des Reziprozitätsgesetzes
der quadratischen Reste her\footnote{\gesperrt{Dirichlet-Dedekind}, Zahlentheorie, 4. Aufl. (S. 636, Anm.).}.

Nach Bd.\ II, \S\,167, 3. ist die notwendige und hinreichende
Bedingung dafür, daß das in der Primzahl $p$ aufgehende Primideal
 $\mathfrak{p}$ (in irgend einem Körper $\Omega$) vom ersten Grade sei, die,
daß für \gesperrt{jede} ganze Zahl $\omega$ des Körpers $\Omega$
\weq{1}{\omega^p \equiv \omega \ (\mathrm{mod}\ \mathfrak{p}),}
%CONT>>
